
\section{Bremsstrahlung of an electron on an electron\\in the ultrarelativistic case}

\SourcePage{472}{477}

The bremsstrahlung of an electron on an electron is depicted by eight Feynman diagrams: four diagrams

% [Feynman diagrams (97.1a): four direct bremsstrahlung diagrams with initial momenta p₁, p₂ and final momenta p'₁, p'₂, photon k emitted from various electron lines]

\vspace{1ex}
\noindent and four ``exchange'' diagrams, obtained from the depicted ones by the interchange of $p_1'$ and $p_2'$. We present here the results of the calculations for the ultrarelativistic case (\emph{G.~Altarelli, F.~Buccella}, 1964; \emph{V.~N.~Baier, V.~S.~Fadin, V.~A.~Khoze}, 1966)\footnote{The corresponding calculations can be found in the book by Baier, Katkov and Fadin indicated on p.~454.}.

In the laboratory frame (the rest frame of one of the initial electrons---say, the second), the cross section of radiation integrated over the photon directions can be represented as a sum $d\sigma = d\sigma^{(1)} + d\sigma^{(2)}$, where
\begin{equation}
d\sigma^{(1)} = 4\alpha r_\mathrm{e}^2\,\frac{d\omega}{\omega}\,\frac{\varepsilon-\omega}{\varepsilon}\!\left[\frac{\varepsilon}{\varepsilon-\omega} + \frac{\varepsilon-\omega}{\varepsilon} - \frac{2}{3}\right]\!\left(\ln\frac{2\varepsilon(\varepsilon-\omega)}{m\omega} - \frac{1}{2}\right);
\label{eq:97.2}
\end{equation}

\SourcePage{473}{478}

\begin{equation}
d\sigma^{(2)} = \frac{2}{3}\,\alpha r_\mathrm{e}^2\,\frac{m\,d\omega}{\omega^2}\!\left\{\!\left(4 - \frac{m}{\omega} + \frac{m^2}{4\omega^2}\right)\ln\frac{2\varepsilon}{m} - 2 + \frac{2m}{\omega} - \frac{5m^2}{8\omega^2}\right\}\; \omega \gg \frac{m}{2};
\label{eq:97.3}
\end{equation}
\begin{equation}
\begin{split}
d\sigma^{(2)} &= \frac{2}{3}\,\alpha r_\mathrm{e}^2\,\frac{d\omega}{\omega}\biggl\{8\!\left(1 - \frac{m}{\omega} + \frac{\omega^2}{m^2}\right)\ln\frac{\varepsilon}{m} - \left(1 - \frac{2\omega}{m}\right)\ln\!\left(1 - \frac{2\omega}{m}\right) \times {} \\
&\quad \times \left[\frac{m^3}{4\omega^3} - \frac{m^2}{2\omega^2} + \frac{3m}{\omega} - \frac{4\omega}{m}\right] - \frac{m^2}{2\omega^2} + \frac{3m}{\omega} - 2 + \frac{m^2}{2\omega^2} - \frac{4\omega^2}{m^2}\biggr\}\; \omega \leqslant \frac{m}{2};
\end{split}
\label{eq:97.4}
\end{equation}
($\varepsilon$ is the initial energy of the first electron).

The accuracy of these formulas is up to terms of relative order $m/\varepsilon$. With this accuracy it turns out that the contributions to the cross section from the various diagrams do not interfere with each other, and in this sense $d\sigma^{(1)}$ and $d\sigma^{(2)}$ correspond to radiation by each of the two electrons---respectively the fast electron and the recoil electron (diagrams (97.1a) and (97.1b)).

The exchange-type diagrams give the same contribution to the cross section as the ``direct'' diagrams. By virtue of the identity of the electrons the total contribution of the direct and exchange diagrams must be divided by 2; therefore one can consider only the contribution of the direct diagrams and not take into account the identity of the particles. For the collision of an electron with a positron, in place of exchange diagrams there appear annihilation diagrams. Their contribution, however, turns out to be of relative order $m/\varepsilon$, i.e.\ negligible. Therefore, with the indicated accuracy, the bremsstrahlung cross sections in electron-electron and electron-positron collisions are the same.

When $\omega \gg m$ the ratio
\[
\frac{d\sigma^{(2)}}{d\sigma^{(1)}} \sim \frac{m}{\omega} \ll 1,
\]
i.e.\ the radiation by the recoil electron is small compared with the radiation by the fast electron (when this ratio reaches order $m/\varepsilon$, formula (97.3), of course, loses its meaning). On the contrary, when $\omega \ll m$ both parts of the cross section are nearly equal:
\begin{equation}
d\sigma^{(1)} = \frac{16}{3}\,\alpha r_\mathrm{e}^2\,\frac{d\omega}{\omega}\ln\frac{2\varepsilon^2}{m\omega}, \quad d\sigma^{(2)} = \frac{16}{3}\,\alpha r_\mathrm{e}^2\,\frac{d\omega}{\omega}\ln\frac{\varepsilon}{\omega}, \quad \omega \ll m.
\label{eq:97.5}
\end{equation}

For the validity of formulas (97.2)--(97.5) it is necessary that at least one of the electrons after the radiation remain ultrarelativistic. In other words, the photon frequency must be sufficiently far from the hard boundary of the spectrum, i.e.\ from the maximum frequency $\omega_\text{max}$ that can be radiated. The final energy of the electrons will be minimal, and the photon energy maximal, when both electrons move after the radiation in the direction of the photon and have equal velocities. Then from the conservation laws we have
\[
\varepsilon + m = \omega_\text{max} + 2\varepsilon', \quad |\mathbf{p}| = \omega_\text{max} + 2|\mathbf{p}'|.
\]

\SourcePage{474}{479}

Eliminating from this $\varepsilon'$ and $\mathbf{p}'$, we obtain
\[
(\varepsilon + m - \omega_\text{max})^2 - (|\mathbf{p}| - \omega_\text{max})^2 = 4m^2,
\]
whence
\begin{equation}
\omega_\text{max} = \frac{m(\varepsilon - m)}{m + \varepsilon - |\mathbf{p}|}.
\label{eq:97.6}
\end{equation}

When $\varepsilon \gg m$ we have $\omega_\text{max} \approx \varepsilon$. Thus, formulas (97.2)--(97.4) are valid under the condition
\begin{equation}
\omega_\text{max} - \omega \sim \varepsilon - \omega \gg m.
\label{eq:97.7}
\end{equation}

The cross section of radiation by the fast electron (97.2) coincides exactly with the cross section of radiation of an electron on a nucleus with $Z = 1$ (formula (93.17)). This coincidence is not accidental, and its causes are elucidated by analysis of the role of recoil in the radiation process.

In deriving formula (93.17) we neglected the recoil of the stationary particle (the nucleus)---replaced it by a constant external field. This reduces to the neglect of the time component of the 4-vector of momentum transfer $q = p' - p + k$ (the recoil energy). Let us show that in the ultrarelativistic case such neglect is permissible for radiation by an electron not only on a nucleus but also on an electron.

Let us write $q^2$ in the form
\begin{equation}
-q^2 = -(\varepsilon' + \omega - \varepsilon)^2 + (p_\parallel' + \omega - p_\parallel)^2 + (\mathbf{p}_\perp' - \mathbf{p}_\perp)^2,
\label{eq:97.8}
\end{equation}
where the subscripts denote the components of the vectors $\mathbf{p}'$ and $\mathbf{p}$ (the initial and final momenta of the electron), parallel and perpendicular to the direction of the photon $\mathbf{k}$. In the ultrarelativistic case the angles $\theta$ and $\theta'$ (between $\mathbf{k}$ and respectively $\mathbf{p}$ and $\mathbf{p}'$) are small: $\theta \lesssim m/\varepsilon$, $\theta' \lesssim m/\varepsilon'$. Therefore
\begin{equation}
|\mathbf{p}_\perp| \sim |\mathbf{p}|\theta \sim m, \quad p_\parallel \approx |\mathbf{p}| - \frac{p_\perp^2}{2|\mathbf{p}|} \approx \varepsilon - \frac{m^2}{2\varepsilon} - \frac{\mathbf{p}_\perp^2}{2\varepsilon},
\label{eq:97.9}
\end{equation}
and analogously for $\mathbf{p}_\perp'$, $p_\parallel'$.

Without account of recoil we have $\varepsilon' + \omega - \varepsilon = 0$, the difference $p_\parallel' + \omega - p_\parallel \sim m^2/\varepsilon$, so that
\begin{equation}
-q^2 \approx (\mathbf{p}_\perp' - \mathbf{p}_\perp)^2 \sim m^2.
\label{eq:97.10}
\end{equation}

The recoil energy (on an electron):
\begin{equation}
q_0 = \varepsilon' + \omega - \varepsilon \sim \mathbf{q}^2/(2m) \sim m.
\label{eq:97.11}
\end{equation}
The change of $\mathbf{p}_\perp'$ due to the change of $\varepsilon'$ can be neglected. Therefore the first two terms in (97.8) give the change of $q^2$ upon account of

\SourcePage{475}{480}

recoil; let us denote it by $\Delta q^2$. Using (97.9), we obtain
\[
\Delta q^2 \approx (\varepsilon' + \omega - \varepsilon)\!\left(-\frac{m^2}{\varepsilon'} - \frac{\mathbf{p}_\perp'^2}{\varepsilon'} + \frac{m^2}{\varepsilon} + \frac{\mathbf{p}_\perp^2}{\varepsilon}\right) \sim m^2\cdot\frac{m}{\varepsilon}.
\]
Comparing with (97.10), we see that $\Delta q^2 \ll |q^2|$, which justifies the neglect of recoil\footnote{This conclusion is, of course, all the more valid for radiation of an electron on a nucleus, for which the recoil energy $q_0 \approx \mathbf{q}^2/(2M) \sim m^2/M$, where $M$ is the mass of the nucleus.}.

The fact that the fast particle radiates into a narrow cone (with opening angle $\sim m/\varepsilon$) in the direction of its motion permits obtaining the radiation cross section in the centre-of-mass frame by a simple recalculation of the cross section (97.2) from the laboratory frame\footnote{In the general case such a recalculation is impossible, since the contribution to the spectrum in a given frequency interval $d\omega$ arises from photons emitted in essentially different directions.}.

In the centre-of-mass frame both electrons radiate equally, each in the direction of its motion (this circumstance visually explains the reason for the absence of interference between the radiation of both particles). The energy of the ultrarelativistic electron in the centre-of-mass frame is related to its energy $\varepsilon$ in the laboratory frame by $2E^2 = m\varepsilon$, and the frequencies $\Omega$ and $\omega$ of the photon in these frames by $\omega/\varepsilon = \Omega/E$ (these equalities are easy to obtain by comparing the values of the invariants $(p_1 p_2)$ and $(p_1 k)$ in both frames). Therefore for the cross section of radiation by each of the electrons in the centre-of-mass frame we find
\begin{equation}
\begin{split}
d\sigma^{(1)} = d\sigma^{(2)} &= 4\alpha r_\mathrm{e}^2\,\frac{d\Omega}{\Omega}\,\frac{E-\Omega}{E}\!\left[\frac{E}{E-\Omega} + \frac{E-\Omega}{E} - \frac{2}{3}\right] \times {} \\
&\quad \times \left(\ln\frac{4E^2(E-\Omega)}{m^2\Omega} - \frac{1}{2}\right).
\end{split}
\label{eq:97.12}
\end{equation}

For the applicability of (97.12) it is also necessary that the photon frequency not be close to the boundary of the spectrum. For an ultrarelativistic particle the transformation indicated above directly gives from $\omega_\text{max} \approx \varepsilon$
\begin{equation}
\Omega_\text{max} \sim \omega_\text{max} E/\varepsilon \approx E.
\label{eq:97.13}
\end{equation}
Thus, in the centre-of-mass frame the electrons can radiate only half their total energy $2E$. A direct computation of $\Omega_\text{max}$ is easily carried out by noting that after radiation of such a photon the electrons will move (in the same frame) with equal velocities in the direction opposite to the direction of the photon. We have
\[
2E = 2E' + \Omega_\text{max}, \quad 2|\mathbf{p}'| = \Omega_\text{max},
\]

\SourcePage{476}{481}

whence
\begin{equation}
\Omega_\text{max} = \frac{\mathbf{p}^2}{m} = E - \frac{m^2}{E},
\label{eq:97.14}
\end{equation}
and in the ultrarelativistic case we again obtain (97.13). Thus, formula (97.12) is applicable under the condition
\begin{equation}
\Omega_\text{max} - \Omega \sim E - \Omega \gg m.
\label{eq:97.15}
\end{equation}

Let us now give formulas for radiation in the centre-of-mass frame in the opposite limiting case, near the boundary of the spectrum, when\footnote{The result obtained in the Born approximation is, of course, valid only so long as the relative velocity of the final electrons is large compared with $\alpha$. In the opposite case one should take into account the interaction of the particles in the final state.}
\begin{equation}
\Omega_\text{max} - \Omega \ll m.
\label{eq:97.16}
\end{equation}
Since in this case the recoil is very substantial, the results differ from the case of scattering on a fixed centre and turn out to be different for electron-electron and electron-positron scattering (\emph{V.~N.~Baier, V.~S.~Fadin, V.~A.~Khoze}, 1967).

In the case of electron scattering on an electron, besides the squares of diagrams (97.1), a contribution to the radiation cross section near the boundary of the spectrum is also given by the products (interference terms) of direct and exchange diagrams in which one and the same initial particle radiates.

% [Feynman diagram: exchange bremsstrahlung diagram with photon k emitted from initial electron p₁, with final states p'₂ and p'₁ (exchanged)]

This is connected with the fact that near the boundary the final particles have close momenta and there are no reasons for the smallness of the exchange terms. The final answer for the cross section:
\begin{equation}
d\sigma = 2\alpha r_\mathrm{e}^2\,\frac{[E(\Omega_\text{max} - \Omega)]^{1/2}}{m}\,\frac{d\Omega}{\Omega_\text{max}}.
\label{eq:97.17}
\end{equation}

In the scattering of an electron on a positron a logarithmically large contribution to the radiation cross section comes from the squares of annihilation diagrams in which the initial particles radiate:

% [Feynman diagrams (97.18): two annihilation bremsstrahlung diagrams for e⁺e⁻ scattering with photon k]

\SourcePage{477}{482}

With non-logarithmic accuracy the squares of other diagrams are also substantial. The interference terms, however, are small. The final answer:
\begin{equation}
d\sigma = 2\alpha r_\mathrm{e}^2\,\frac{[E(\Omega_\text{max} - \Omega)]^{1/2}}{m}\!\left(\ln\frac{2E}{m} + 1\right)\frac{d\Omega}{\Omega_\text{max}}.
\label{eq:97.19}
\end{equation}
Thus, radiation in electron-positron scattering is logarithmically large compared with radiation in electron-electron scattering.
